\begin{historical}
Ether-era anaesthesia began in 1846 with Morton’s public demonstration in Boston; within months, ether and chloroform spread worldwide. By the early twentieth century, Meyer and Overton independently observed a correlation: anaesthetic potency scaled with lipid solubility across diverse compounds. This supported the idea that consciousness could be turned off by a nonspecific action on neuronal membranes. Yet the correlation cracked under scrutiny: highly lipophilic yet inert molecules failed to anaesthetise, while effective agents deviated from the predicted potency.

Mid-to-late twentieth century work shifted toward specific molecular targets: GABA\textsubscript{A} receptors, NMDA antagonism, two-pore K\textsuperscript{+} channels, and hyperpolarisation-activated cyclic nucleotide–gated (HCN1) currents. Still, no single pathway unified the class.

In prion disease, a different historical thread exposed the opposite failure mode. In 1982, Prusiner proposed prions—a proteinaceous infectious particles—as agents of neurodegeneration. A rare PRNP mutation producing fatal familial insomnia (FFI) was later traced to selective thalamic degeneration, abolishing sleep despite otherwise preserved wakeful function. An Italian pedigree provided the defining clinical arc: onset with fragmented sleep, inexorable insomnia, autonomic failure, cognitive collapse, and death within months. Where anaesthesia induced obliviousness, FFI prevented it.

\end{historical}
